\section{Codebook for the Expert Interviews}
\label{annex:expert-interview-codebook}

This is the codebook used for the three expert interviews. The interviews do two kinds of work in this thesis.
They describe how a remote group stays with a joint decision, and they are the practical basis
for requirements and for several user stories of the AI moderator.

The Coding-Scheme has two levels: four categories and fifteen codes.
A passage is always coded with a code, never with a category alone.

\begin{itemize}
  \item \textbf{Decision conditions (D)} describe what has to be true around a decision, including
    work before and after the meeting.
  \item \textbf{Meeting dynamics (M)} describe what puts goal orientation or decision quality at
    risk while the group is together, including what is specific to the remote setting.
  \item \textbf{Facilitation practice (F)} describes what the experts do about that.
  \item \textbf{AI co-facilitator (A)} records where a system would help, how it should enter the
    meeting, and what it should not take over.
\end{itemize}

The clusters in \autoref{annex:expected-thematic-clusters} were a sensitizing frame for the guide.
They are not this coding scheme.

\subsection*{Coding rules}

\begin{itemize}
  \item \textbf{Unit.}
    Code a coherent claim by the interviewee. A claim may run across consecutive rows of the same
    speaker. Interviewer turns, consent, and scheduling are not coded.
  \item \textbf{Multiple coding.}
    One claim can take more than one code when it states both a problem and a response, or both a
    useful function and a limit.
  \item \textbf{Category cut.}
    D and M may both apply to one claim.
    D is what the meeting depends on or cannot remove: the specification of the decision,
    preparation and the handoff, and structure.
    M is what that does while people are together: loss of the question, who speaks or stays
    silent, and a cue that cannot be read.
    The clock is not the cut.
    A structural limit felt inside the meeting is still D3.
    A fear culture that shapes who speaks is M2, and it is also D3 when the speaker says one
    session cannot change it.
  \item \textbf{Explicitness.}
    Code what is said or clearly implied. Do not invent a motive.
    A method name is not a code. Code the job the method does, such as structuring the process,
    making the state visible, or bringing in a voice.
  \item \textbf{Grain.}
    Do not open a code for a single technique. Check-in, breakouts, fist-to-five, and a
    prioritization matrix stay under the code for their function.
  \item \textbf{No separate polarity.}
    Problem, practice, valued function, and refusal are carried by the code itself.
\end{itemize}

\begingroup
\footnotesize
\setlength{\tabcolsep}{4pt}
\setlength{\extrarowheight}{0.4pt}
\renewcommand{\arraystretch}{1.05}
\setlength{\LTpre}{0.4em}
\setlength{\LTpost}{0.6em}
\setlength{\emergencystretch}{2em}
\newcommand{\cbgap}{\addlinespace[0.45em]}
\newcommand{\cbcatbreak}{%
\addlinespace[0.15em]\midrule\addlinespace[0.4em]}
\newcommand{\cbcat}[1]{%
\multicolumn{3}{@{}l@{}}{\textbf{#1}} \\*[0.2em]}
\newcommand{\cbrow}[5]{%
  \textbf{#1}\par\vspace{0.1em}#2
  & #3\par\vspace{0.25em}\textit{Rule:} #4
  & \enquote{#5} \\
}
\begin{xltabular}{\textwidth}{@{}
    >{\RaggedRight\arraybackslash}p{3.3cm}
    >{\RaggedRight\arraybackslash}X
    >{\RaggedRight\arraybackslash}X
  @{}}
  \caption{Codebook for the expert interviews.
  Anchors are shortened from the transcript and keep the original German wording.}
  \label{tab:expert-interview-codebook}\\[0.7em]
  \toprule
  \textbf{Code} & \textbf{Definition} & \textbf{Anchor example} \\
  \midrule
  \endfirsthead
  \caption[]{Codebook for the expert interviews (continued).}\\[0.7em]
  \toprule
  \textbf{Code} & \textbf{Definition} & \textbf{Anchor example} \\
  \midrule
  \endhead
  \bottomrule
  \endfoot
  \bottomrule
  \endlastfoot
  \cbcat{Decision conditions}
  \cbrow
  {D1}
  {Decision scope, process and rights}
  {What is being decided, who may decide, and by which procedure. Includes authority,
    majority, consensus, and a consultative single decision, and whether participation is real
    or only a formal nod. Also the difference between a reversible and an irreversible decision,
  and that not deciding is itself a decision.}
  {The object, the rights, or the procedure of the decision. Running a method inside the
  meeting is F, not D1. The effect of that procedure on who speaks or who loses is M2.}
  {Die meisten kennen in Unternehmen 2 Arten von Entscheidungsverfahren, Chef entscheidet
  [\ldots] und das andere Entscheidungsverfahren ist Mehrheit.}
  \cbgap
  \cbrow
  {D2}
  {Before and after the meeting}
  {Work outside the live discussion: preparation, a shared information base, unprepared
    participants, documentation, and the handoff into next steps. Includes who does what by
  when, and whether the result is checked afterwards.}
  {Conditions outside the session. A summary of the current discussion is F2.}
  {das allergrößte Manko wirklich bei Remote-Meetings, dass Leute in die Meetings reinkommen
  und nicht vorbereitet sind.}
  \cbgap
  \cbrow
  {D3}
  {Structural limits of the meeting}
  {The meeting as a symptom of something the meeting itself cannot fix: missing autonomy,
    the wrong people or too many people, a meeting used only to broadcast information, or a
  decision postponed until the next regular slot.}
  {The speaker says that better moderation, or an AI, would not remove the underlying
  problem. If the same claim describes silence or dissent in the room, add M2.
  D3 is only the claim that the meeting cannot remove the cause.}
  {du kriegst diese strukturellen Probleme nicht durch besseres Meeting.}
  \cbcatbreak
  \cbcat{Meeting dynamics}
  \cbrow
  {M1}
  {Focus, time, and the groan zone}
  {The group leaves the decision: side tracks, circular talk, a decision that is not ripe,
    or time pressure that rushes or avoids the uncomfortable middle. Also discovering late that
  the group has been working on the wrong problem.}
  {The problem or the critical moment. The remedy is an F code.
  A decision that is not ripe because the question or the preparation is missing is also
  D1 or D2. M1 is the drift, the side track, or the premature close.}
  {ständig [\ldots] Nebenschauplätze aufgemacht werden. Vielleicht weil die Entscheidung
  einfach noch nicht reif ist.}
  \cbgap
  \cbrow
  {M2}
  {Voices, hierarchy, and dissent}
  {Unequal speaking, dominance, hierarchy, and whether criticism can be said. Includes trust
    and fear culture as conditions for dissent, and silence of someone whose stake matters.
  Silence of someone with nothing to add is not this code.}
  {The social dynamic. A concrete invitation is F3.
  Hierarchy as a culture the facilitator cannot change is also D3.
  The procedure that produces the silence is also D1.}
  {Diese Kultur, dass Widersprüche immer etwas Positives sind [\ldots], leben ganz wenige
  Organisationen.}
  \cbgap
  \cbrow
  {M3}
  {Remote conditions and ambiguous signals}
  {What changes because the meeting is remote: cameras off, parallel work, missing nonverbal
    and cultural cues, no informal clarification in a break, and a reluctant yes that is hard
    to tell from agreement. Also tools and a shared visual surface as a condition of a complex
  decision online.}
  {The remote setting or an unreliable cue. Use A4 as well when the point is that an AI would
  lack that cue.}
  {Wenn ich online eine Pause mache, sind die Leute erst mal Disconnected und da können sich
  vielleicht so kleinere Punkte dann nicht noch klären.}
  \cbcatbreak
  \cbcat{Facilitation practice}
  \cbrow
  {F1}
  {Process structure}
  {A deliberate shape for the meeting: agenda, check-in, small groups, opening before
    closing, separating clarification from debate, and switching to the meta level near the end
  (how the group will decide, instead of reopening the content).}
  {The process design. A single paraphrase is F2.}
  {Du machst ein Check-in, dass du am Anfang schon mal alle Stimmen im Raum hast.}
  \cbgap
  \cbrow
  {F2}
  {Sensemaking and a visible state}
  {Condensing or externalizing where the group stands: summary, clustering, paraphrase,
    options with consequences, and a display of positions that everyone can see at once, such
  as a scale, a poll, or a matrix.}
  {An act that makes the discussion state visible. A pure time warning is A1 when it is
  proposed for the AI, and F4 when it is a judgment about intervening.}
  {Sensemaking miteinander zu machen [\ldots] diese Ergebnisse nochmal in eine Clusterung
  bringst [\ldots] das und das und das die Cluster sind.}
  \cbgap
  \cbrow
  {F3}
  {Bringing in voices and information}
  {A move that gets a missing person or a missing fact into the discussion: a named
    invitation, hearing lower hierarchy before leadership, writing as well as speaking, and
    holding back evaluation while positions are still being collected. Includes asking what an
  absent role would say.}
  {The move, not the mere fact that voices are unequal (M2). Role matters: inviting everyone
  equally is not required by this code.}
  {Die Praktikantin, die nichts sagt [\ldots], die werde ich nicht ansprechen, aber den
  Co-founder, der länger nichts gesagt hat [\ldots], den werde ich schon mal fragen.}
  \cbgap
  \cbrow
  {F4}
  {Intervention judgment}
  {When to step in, when to wait, and whom to address. Includes holding silence, the risk of
    being too early or too late, and stopping when a non-answer is itself information. Also
  face and tone: an invitation that would expose someone.}
  {The judgment about timing, addressee, or restraint. The content of the move may also
  receive F1, F2, or F3.}
  {manchmal schweige ich auch einfach und warte [\ldots]. wenn nach 7~s keiner [\ldots], dann
  weiß ich [\ldots], da muss ich jetzt nochmal [\ldots] anders fragen.}
  \cbgap
  \cbrow
  {F5}
  {Neutrality and ownership}
  {The facilitator structures the conversation and does not make the decision. The group, or
    a named decision owner, has to take the result as their own. A decision cannot be delegated
  to the facilitator.}
  {The human role. The same boundary stated for an AI is A3.}
  {die Entscheidung wirklich zu treffen [\ldots] ist nicht die Aufgabe von Facilitator,
  sondern die Aufgabe der Führungskraft.}
  \cbcatbreak
  \cbcat{AI co-facilitator}
  \cbrow
  {A1}
  {Promising functions}
  {A concrete use the speaker treats as helpful: timekeeping and agenda signals, fast
    summaries and clustering, walking a group through a decision (question, options, doing
    nothing, pros and cons), pointing at a decision the group has not noticed, a bias cue,
  documentation into existing tools, or a persona for a missing stakeholder.}
  {A valued function. If the same claim limits when it may happen, add A2 or A3.}
  {wenn sie es schafft, schneller diese Zusammenfassungen zu machen [\ldots] aus komplexen
  Daten schneller Sinn zu machen.}
  \cbgap
  \cbrow
  {A2}
  {Role, trigger, and acceptance}
  {How the system should enter the meeting and whether people would accept it. Includes
    on-demand versus proactive speech, a prior agreement of the group, a human who triggers
    the system, and presence: a chat is talked over, and people may not yield the floor to a
  machine.}
  {The interaction role, not the function itself (A1).}
  {Muss aber eingeführt werden. Und muss ein Agreement der ganzen Gruppe haben, dass die das
  machen kann.}
  \cbgap
  \cbrow
  {A3}
  {Risks and non-goals}
  {Harm, refusal, or a boundary: surveillance and sentiment analysis, data protection,
    exposing quiet people, a false read, the system deciding or imposing a structure nobody
  asked for, and efficiency that crowds out human talk.}
  {A limit or a harm. A function that is useful under a condition is A1 plus A2 or A3.}
  {das kann ganz schnell in Kontrolle übergehen. [\ldots] die letzten zehn Meetings hast du
  dich kein einziges Mal gemeldet.}
  \cbgap
  \cbrow
  {A4}
  {Context the system does not have}
  {Knowledge a facilitator uses that sits outside the transcript: roles and stakes, the
    decision scope, process and rights, and earlier meetings, private circumstances, nonverbal cues, and the culture of
  the group.}
  {A missing context as a design constraint. The human difficulty of reading the remote room
  stays M3 unless the AI is the point.}
  {hier bei jemandem ist gerade die Tochter krank. [\ldots] das weiß die KI aber nicht.}
\end{xltabular}
\endgroup
